%Version 3.1 December 2024
% See section 11 of the User Manual for version history
%
%%%%%%%%%%%%%%%%%%%%%%%%%%%%%%%%%%%%%%%%%%%%%%%%%%%%%%%%%%%%%%%%%%%%%%
%%                                                                 %%
%% Please do not use \input{...} to include other tex files.       %%
%% Submit your LaTeX manuscript as one .tex document.              %%
%%                                                                 %%
%% All additional figures and files should be attached             %%
%% separately and not embedded in the \TeX\ document itself.       %%
%%                                                                 %%
%%%%%%%%%%%%%%%%%%%%%%%%%%%%%%%%%%%%%%%%%%%%%%%%%%%%%%%%%%%%%%%%%%%%%

%%\documentclass[referee,sn-basic]{sn-jnl}% referee option is meant for double line spacing

%%=======================================================%%
%% to print line numbers in the margin use lineno option %%
%%=======================================================%%

%%\documentclass[lineno,pdflatex,sn-basic]{sn-jnl}% Basic Springer Nature Reference Style/Chemistry Reference Style

%%=========================================================================================%%
%% the documentclass is set to pdflatex as default. You can delete it if not appropriate.  %%
%%=========================================================================================%%

%%\documentclass[sn-basic]{sn-jnl}% Basic Springer Nature Reference Style/Chemistry Reference Style

%%Note: the following reference styles support Namedate and Numbered referencing. By default the style follows the most common style. To switch between the options you can add or remove “Numbered” in the optional parenthesis.
%%The option is available for: sn-basic.bst, sn-chicago.bst%

%%\documentclass[pdflatex,sn-nature]{sn-jnl}% Style for submissions to Nature Portfolio journals
%%\documentclass[pdflatex,sn-basic]{sn-jnl}% Basic Springer Nature Reference Style/Chemistry Reference Style
\documentclass[pdflatex,sn-mathphys-num]{sn-jnl}% Math and Physical Sciences Numbered Reference Style
%%\documentclass[pdflatex,sn-mathphys-ay]{sn-jnl}% Math and Physical Sciences Author Year Reference Style
%%\documentclass[pdflatex,sn-aps]{sn-jnl}% American Physical Society (APS) Reference Style
%%\documentclass[pdflatex,sn-vancouver-num]{sn-jnl}% Vancouver Numbered Reference Style
%%\documentclass[pdflatex,sn-vancouver-ay]{sn-jnl}% Vancouver Author Year Reference Style
%%\documentclass[pdflatex,sn-apa]{sn-jnl}% APA Reference Style
%%\documentclass[pdflatex,sn-chicago]{sn-jnl}% Chicago-based Humanities Reference Style

%%%% Standard Packages
%%<additional latex packages if required can be included here>

\usepackage{graphicx}%
\usepackage{multirow}%
\usepackage{amsmath,amssymb,amsfonts}%
\usepackage{amsthm}%
\usepackage{mathrsfs}%
\usepackage[title]{appendix}%
\usepackage{xcolor}%
\usepackage{textcomp}%
\usepackage{manyfoot}%
\usepackage{booktabs}%
\usepackage{algorithm}%
\usepackage{algorithmicx}%
\usepackage{algpseudocode}%
\usepackage{listings}%
%%%%

%%%%%=============================================================================%%%%
%%%%  Remarks: This template is provided to aid authors with the preparation
%%%%  of original research articles intended for submission to journals published
%%%%  by Springer Nature. The guidance has been prepared in partnership with
%%%%  production teams to conform to Springer Nature technical requirements.
%%%%  Editorial and presentation requirements differ among journal portfolios and
%%%%  research disciplines. You may find sections in this template are irrelevant
%%%%  to your work and are empowered to omit any such section if allowed by the
%%%%  journal you intend to submit to. The submission guidelines and policies
%%%%  of the journal take precedence. A detailed User Manual is available in the
%%%%  template package for technical guidance.
%%%%%=============================================================================%%%%

%% as per the requirement new theorem styles can be included as shown below
\theoremstyle{thmstyleone}%
\newtheorem{theorem}{Theorem}%  meant for continuous numbers
%%\newtheorem{theorem}{Theorem}[section]% meant for sectionwise numbers
%% optional argument [theorem] produces theorem numbering sequence instead of independent numbers for Proposition
\newtheorem{proposition}[theorem]{Proposition}%
%%\newtheorem{proposition}{Proposition}% to get separate numbers for theorem and proposition etc.

\theoremstyle{thmstyletwo}%
\newtheorem{example}{Example}%
\newtheorem{remark}{Remark}%

\theoremstyle{thmstylethree}%
\newtheorem{definition}{Definition}%

\raggedbottom
%%\unnumbered% uncomment this for unnumbered level heads

\usepackage{amsmath}
\usepackage{mathpartir}

\newcommand{\alg}{\Rightarrow}

\begin{document}

\title{Coercive Subtyping via Logic Programming for Painless Mendler-style Algebraic Programming}

\author*[1]{\fnm{Henry} \sur{Blanchette}}
\email{blancheh@umd.edu}

\author*[2]{Aaron Stump}
\email{stumpaa@bc.edu}

\affil[1]{\orgname{University of Maryland}}
\affil[2]{\orgname{Boston College}}

% \begin{CCSXML}
%   <ccs2012>
%   <concept>
%   <concept_id>10011007.10011006.10011008.10011009.10011012</concept_id>
%   <concept_desc>Software and its engineering~Functional languages</concept_desc>
%   <concept_significance>500</concept_significance>
%   </concept>
%   <concept>
%   <concept_id>10011007.10011006.10011008.10011024.10011033</concept_id>
%   <concept_desc>Software and its engineering~Recursion</concept_desc>
%   <concept_significance>500</concept_significance>
%   </concept>
%   </ccs2012>
% \end{CCSXML}

% \ccsdesc[500]{Software and its engineering~Functional languages}
% \ccsdesc[300]{Software and its engineering~Recursion}

\keywords{strong functional programming, type-based termination, Mendler-style recursion}

\abstract{
  This paper presents a system for more convenient algebraic programming, a form of total functional programming where datatypes are defined as least fixed points of covariant signature functors, and recursive functions are defined as algebras. This approach requires the use of certain combinators to roll and unroll fixed points, and to convert algebras to functions which can be applied to recurse over data. These combinators are a burden for programming in this style. To alleviate that burden, in this paper we propose to insert them automatically, using coercive subtyping. The type checker generates subtyping constraints, which are then solved with the help of a custom logic programming engined called Chronolog. The subtyping axioms are given to Chronolog as rules, which uses them to solve the subtyping constraints. To avoid infinite applications of rules for unrolling signature functors, Chronolog provides a mechanism for suspending and resuming certain forms of goals. The subtyping axioms presented to Chronolog are an algorithmic version of the usual declarative rules for algebraic programming. We prove equivalence of the declarative and algorithmic rules, in Agda. We present numerous examples of algebraic programs written in our system.
}

\maketitle

\section{Introduction} % Aaron

% Algebraic programming as an example of total functional programming (cite Turner)
% Mendler-style algebra
% Burden of inserting roll/unroll and cata
% We will reduce this burden using coercive subtyping to infer the coercions

\section{Motivating example} % Henry

% pick nice example
%
%   subtraction uses unroll, but not roll or cata
%   length uses roll, but not unroll or cata
%   write an outer function that subtracts some quantity from the length of a list
%
% here it is in Haskell
% here it is in algebraic style with coercions
%
% Can we use the algebraic style with lighter syntax?  This would make it more
% feasible in practice.

\section{Syntax} % Henry (maybe Aaron helping with prose)

\section{Constraint-Generating Typing} % Henry

\begin{itemize}
  \item $\mathscr{D}$ is the set of datatype definitions, which are pairs of $\langle D , V \rangle$ where $D$ is an identifier and $V$ is a set of constructor signatures.
  \item The meta-function $\textsf{freshenFunctorParam}(A)$ replaces all instances of the functor parameter type in $A$ with a fresh type variable.
  \item The meta-function $\textsf{typeOfCoercion}(c)$ calculates the pair of the domain and codomain of a coercion.
\end{itemize}

\newcommand{\sLet}{\text{let}~}
\newcommand{\sIn}{~\text{in}~}
\newcommand{\sCoerce}{\text{coerce}~}

\begin{figure}[H]
  \centering

  \begin{mathpar}

    \inferrule[InferVar]{
      (x : A) \in \Gamma
    }{
      \Gamma \vdash x : A ~|~ \varnothing
    }

    \inferrule[InferCon]{
      \langle D , V \rangle \in \mathscr{D} \quad
      \langle c , \langle \dots A'_i, \dots \rangle \rangle \in V
      \\\\
      (\forall i)~ \Gamma \vdash a_i : A_i ~|~ \mathcal{C}_{a_i}
    }{
      \Gamma \vdash D.c(\dots a_i, \dots) ~|~ \bigcup_{i} (\mathcal{C}_{a_i} \cup \{ A_i <: \textsf{freshenFunctorParam}(A'_i) \})
    }

    \inferrule[InferLam]{
      \Gamma, (x : A) \vdash b : B ~|~ \mathcal{C}_b
    }{
      \Gamma \vdash \lambda (x : A) . b : A \to B ~|~ \mathcal{C}_b
    }

    \inferrule[InferApp]{
      A', B ~\text{are fresh type variables}
      \\\\
      \Gamma \vdash f : F ~|~ \mathcal{C}_f \quad
      \Gamma \vdash a : A ~|~ \mathcal{C}_a
    }{
      \Gamma \vdash (f ~ a) : B ~|~ \mathcal{C}_f \{ F <: A' \to B \} \cup \mathcal{C}_a \cup \{ A <: A' \}
    }

    \inferrule[InferGamma]{
      \langle D , V \rangle \in \mathscr{D} \\\\
      (\forall i)~ \langle c_i , \langle \dots A'_{ij}, \dots \rangle \rangle \in V
      \\\\
      (\forall i)~ \Gamma, \dots x_{ij} : A'_{ij} \dots \vdash b_i : B_i ~|~ \mathcal{C}_{b_i}
    }{
      \Gamma \vdash \gamma(D, B) a . \{ \dots ~ c_i(\dots x_{ij}, \dots) \mapsto b_i | ~ \dots \} ~|~ \bigcup_i (\mathcal{C}_{b_i} \cup \{ B_i <: B \})
    }

    \inferrule[InferOmega]{
      \Gamma, (R : \ast), (f : R \to A), (x : D ~ R) \vdash b : B ~|~ \mathcal{C}_b
    }{
      \Gamma \vdash \omega(D, B, R, f, x) . b : D \alg B ~|~ \mathcal{C}_b
    }

    \inferrule[InferLet]{
      \Gamma \vdash a : A ~|~ \mathcal{C}_a \quad
      \Gamma, (x : A') \vdash b : B ~|~ \mathcal{C}_b
    }{
      \Gamma \vdash \sLet x : A' = a \sIn b : B ~|~ \mathcal{C}_a \cup \{ A <: A' \} \cup \mathcal{C}_a
    }

    \inferrule[InferCoerce]{
      \langle A', B \rangle := \textsf{typeOfCoercion}(c) \\\\
      \Gamma \vdash a : A ~|~ \mathcal{C}_a
    }{
      \Gamma \vdash \sCoerce c ~ a : B ~|~ \mathcal{C}_a \cup \{ A <: A' \}
    }

  \end{mathpar}

  \caption{Subtype constraint generation rules.}
  \label{fig:subtype-constraint-generation-rules}
\end{figure}

\section{Subtyping} % Aaron

\subsection{Declarative}

\begin{figure}[H]
  \centering

  \begin{mathpar}

    \inferrule[SubData]
    {
      % D is a datatype
      \langle D, \_ \rangle \in \mathscr{D}
    }
    {
      D <: D
    }

    \inferrule[Omega]
    {
      % R is an omega var
      \text{TODO: require that R is an omega var}
    }
    {
      R <: R
    }

    \inferrule[Arr]
    {
      A' <: A \quad
      B <: B'
    }
    {
      A \to B <: A' \to B'
    }

    \inferrule[Cov]{
      A <: A'
    }{
      D ~ A <: D ~ A'
    }

    \inferrule[Alg]{
      A <: A'
    }{
      D \alg A <: D \alg A'
    }

    \inferrule[Roll]{ }{
      D ~ D <: D
    }

    \inferrule[Unroll]{ }{
      D <: D ~ D
    }

    \inferrule[Cata]{ }{
      D \alg A <: D \to A
    }

    \inferrule[Trans]{
      A <: A' \quad
      A' < A''
    }{
      A <: A''
    }
  \end{mathpar}

  \caption{Declarative subtyping rules.}
  \label{fig:declarative-subtyping-rules}
\end{figure}

\subsection{Algorithmic}

\begin{figure}[H]
  \centering

  \begin{mathpar}

    \inferrule[DecSubData]
    {
      % D is a datatype
      \langle D, \_ \rangle \in \mathscr{D}
    }
    {
      D <: D
    }

    \inferrule[DecOmega]
    {
      % R is an omega var
      \text{TODO: require that R is an omega var}
    }
    {
      R <: R
    }

    \inferrule[DecArr]
    {
      A' <: A \quad
      B <: B'
    }
    {
      A \to B <: A' \to B'
    }

    \inferrule[DecCov]{
      A <: A'
    }{
      D ~ A <: D ~ A'
    }

    \inferrule[DecAlg]{
      A <: A'
    }{
      D \alg A <: D \alg A'
    }

    \inferrule[DecRoll]{
      A <: D
    }{
      D ~ A <: D
    }

    \inferrule[DecUnroll]{
      A <: D
    }{
      D <: D ~ A
    }

    \inferrule[DecCata]{
      B <: D \quad
      A <: C
    }{
      D \alg A <: B \to C
    }
  \end{mathpar}

  \caption{Algorithmic subtyping rules.}
  \label{fig:algorithmic-subtyping-rules}
\end{figure}

\subsection{Equivalence of declarative and algorithmic subtyping}

\section{Implementation with Chronolog} % Henry

\section{Returning to the Examples} % Henry

% remind the version with coercions
% show without coercions

\section{Conclusion and future work} % Aaron

% Richer patterns of recursion, as in POPL 2023 paper

\end{document}
